Area ratios are unchanged by an affine coordinate map, so use $A=(0,0),B=(6,0),C=(0,6)$. Then $E=(0,3)$, $F=(3,0)$, $G=(2,2)$, $Y=(3,\dfrac{3}{2})$, $Z=(\dfrac{3}{2},3)$. The area of $GYZ$ is $\dfrac{1}{2}|1\cdot1-(-\dfrac{1}{2})(-\dfrac{1}{2})|=\dfrac{3}{8}$, while $[ABC]=18$. Their ratio is $\dfrac{1}{48}$. Thus the required area is $480\cdot\dfrac{1}{48}=\boxed{10}$.
\begin{asy}
texpreamble("\usepackage{amsmath,amssymb}");
size(240);
defaultpen(fontsize(9pt));
pair A=(0,0),B=(6,0),C=(0,6),E=(0,3),F=(3,0),G=(2,2),Y=(3,1.5),Z=(1.5,3);draw(A--B--C--cycle);draw(B--E);draw(C--F);filldraw(G--Y--Z--cycle,lightblue);label("A",A,SW);label("B",B,SE);label("C",C,N);label("E",E,W);label("F",F,S);label("G",G,SW);label("Y",Y,SE);label("Z",Z,NW);
\end{asy}
